\documentclass[12pt,a4paper]{article}

% Packages
\usepackage[utf8]{inputenc}
\usepackage[T1]{fontenc}
\usepackage[french]{babel}
\usepackage{geometry}
\geometry{left=2.5cm,right=2.5cm,top=2.5cm,bottom=2.5cm}
\usepackage{hyperref}
\hypersetup{
    colorlinks=true,
    linkcolor=blue,
    filecolor=magenta,      
    urlcolor=cyan,
    pdftitle={Documentation d'Architecture - Compilateur Deca},
    pdfauthor={GL 51},
}
\usepackage{listings}
\usepackage{xcolor}
\usepackage{graphicx}
\usepackage{fancyhdr}
\usepackage{tocloft}
\usepackage{enumitem}

% Configuration listings pour Java
\definecolor{javared}{rgb}{0.6,0,0}
\definecolor{javagreen}{rgb}{0.25,0.5,0.35}
\definecolor{javapurple}{rgb}{0.5,0,0.35}
\definecolor{javadocblue}{rgb}{0.25,0.35,0.75}
\definecolor{backcolour}{rgb}{0.95,0.95,0.95}

\lstdefinestyle{javastyle}{
    language=Java,
    backgroundcolor=\color{backcolour},
    commentstyle=\color{javagreen},
    keywordstyle=\color{javapurple}\bfseries,
    numberstyle=\tiny\color{gray},
    stringstyle=\color{javared},
    basicstyle=\ttfamily\footnotesize,
    breakatwhitespace=false,
    breaklines=true,
    captionpos=b,
    keepspaces=true,
    numbers=left,
    numbersep=5pt,
    showspaces=false,
    showstringspaces=false,
    showtabs=false,
    tabsize=2,
    frame=single,
    rulecolor=\color{black}
}

\lstdefinestyle{bashstyle}{
    language=bash,
    backgroundcolor=\color{backcolour},
    basicstyle=\ttfamily\footnotesize,
    breaklines=true,
    frame=single,
    numbers=none,
    commentstyle=\color{javagreen}
}

\lstset{style=javastyle}

% En-têtes et pieds de page
\pagestyle{fancy}
\fancyhf{}
\fancyhead[L]{Documentation d'Architecture}
\fancyhead[R]{GL 51}
\fancyfoot[C]{\thepage}

% Titre
\title{\textbf{Documentation d'Architecture\\Compilateur Deca}}
\author{Équipe GL 51\\Ensimag}
\date{13 janvier 2026}

\begin{document}

\maketitle
\thispagestyle{empty}

\begin{center}
\vspace{2cm}
\large
\textbf{Version 1.0}
\end{center}

\newpage
\tableofcontents
\newpage

\section{Introduction}

\subsection{Contexte et Objectifs}

Ce document décrit l'architecture interne du compilateur Deca. Il s'adresse aux développeurs qui maintiendront ou feront évoluer ce compilateur. Le compilateur Deca traduit des programmes Deca (langage similaire à Java) en assembleur IMA. Ce document explique \textbf{comment} le compilateur fonctionne, pas comment l'utiliser.

\subsection{Qui ? Quoi ? Pourquoi ? Quand ? Pour Qui ?}

\begin{itemize}[itemsep=0.2em]
    \item \textbf{Qui} : L'équipe GL 51, projet Ensimag 2026
    \item \textbf{Quoi} : Un compilateur complet (lexer, parser, analyse contextuelle, génération de code)
    \item \textbf{Pourquoi} : Traduire Deca vers IMA pour exécution sur machine abstraite
    \item \textbf{Quand} : Développé entre décembre 2025 et janvier 2026
    \item \textbf{Pour qui} : Futurs mainteneurs du code, développeurs internes
\end{itemize}

\section{Organisation Générale}

\subsection{Pipeline de Compilation}

Le compilateur suit un pipeline classique en quatre phases :

\begin{enumerate}[itemsep=0.3em]
    \item \textbf{Analyse lexicale} : Découpe le fichier source en tokens (\texttt{DecaLexer.g4})
    \item \textbf{Analyse syntaxique} : Construit l'arbre de syntaxe abstraite (\texttt{DecaParser.g4})
    \item \textbf{Analyse contextuelle} : Vérifie types et portées (méthodes \texttt{verify*()})
    \item \textbf{Génération de code} : Produit l'assembleur IMA (méthodes \texttt{codeGen*()})
\end{enumerate}

\subsection{Orchestration : \texttt{DecacCompiler}}

La classe \texttt{DecacCompiler} coordonne toutes les phases. Elle contient :

\begin{itemize}[itemsep=0.2em]
    \item \texttt{symbolTable} : Table des symboles globale (internement de chaînes)
    \item \texttt{environmentType} : Types prédéfinis (int, float, boolean, void)
    \item \texttt{program} : Représentation IMA accumulant les instructions
    \item \texttt{registerManager} : Gestionnaire d'allocation de registres
\end{itemize}

\textbf{Flux d'exécution} :

\begin{lstlisting}[style=bashstyle]
compile() -> doCompile() -> doLexingAndParsing() 
          -> verifyProgram() -> codeGenProgram()
\end{lstlisting}

Chaque méthode délègue au \texttt{DecacCompiler} pour ajouter des instructions :

\begin{lstlisting}[style=javastyle]
compiler.addInstruction(new LOAD(...));
compiler.addComment("debut boucle");
\end{lstlisting}

\section{Description des Paquetages et Classes Principales}

\subsection{\texttt{fr.ensimag.deca.syntax}}

Contient les grammars ANTLR4 et le code généré.

\begin{itemize}[itemsep=0.2em]
    \item \texttt{DecaLexer.g4} : Règles lexicales (tokens : \texttt{IF}, \texttt{WHILE}, \texttt{IDENT}, etc.)
    \item \texttt{DecaParser.g4} : Règles syntaxiques (production AST)
    \item Classes générées : \texttt{DecaLexer}, \texttt{DecaParser}, listeners
\end{itemize}

\textbf{Dépendance} : ANTLR4 runtime 4.13.2

\subsection{\texttt{fr.ensimag.deca.tree}}

Représente l'arbre de syntaxe abstraite (AST). Chaque nœud hérite de \texttt{Tree}.

\textbf{Hiérarchie principale} :

\begin{lstlisting}[style=bashstyle, numbers=none]
Tree (racine)
├── AbstractProgram → Program
├── AbstractMain → Main, EmptyMain
├── AbstractExpr
│   ├── AbstractBinaryExpr
│   │   ├── AbstractOpArith → Plus, Minus, Multiply, Divide, Modulo
│   │   ├── AbstractOpCmp → Equals, NotEquals, Greater, Lower, etc.
│   │   └── AbstractOpBool → And, Or
│   ├── AbstractUnaryExpr → Not, UnaryMinus
│   ├── AbstractLiteral → IntLiteral, FloatLiteral, BooleanLiteral
│   └── Identifier, ReadInt, ReadFloat
└── AbstractInst
    ├── IfThenElse, While
    ├── Assign, Print, Println
    └── NoOperation
\end{lstlisting}

\textbf{Responsabilités} :

\begin{itemize}[itemsep=0.2em]
    \item Chaque nœud implémente \texttt{verifyExpr()} ou \texttt{verifyInst()} (phase contextuelle)
    \item Chaque nœud implémente \texttt{codeGenExpr()} ou \texttt{codeGenInst()} (phase génération)
    \item \texttt{decompile()} permet de restituer le code source
\end{itemize}

\textbf{Dépendances} :

\begin{itemize}[itemsep=0.2em]
    \item \texttt{fr.ensimag.deca.context} pour types et environnements
    \item \texttt{fr.ensimag.ima.pseudocode} pour instructions IMA
\end{itemize}

\subsection{\texttt{fr.ensimag.deca.context}}

Gère l'analyse contextuelle (vérification de types et portées).

\textbf{Classes clés} :

\begin{itemize}[itemsep=0.2em]
    \item \texttt{Type} : Hiérarchie des types (IntType, FloatType, BooleanType, ClassType, etc.)
    \item \texttt{EnvironmentType} : Environnement global des types
    \item \texttt{EnvironmentExp} : Environnement local pour variables (structure de pile)
    \item \texttt{Definition} : Définitions liées aux symboles (VariableDefinition, MethodDefinition, etc.)
\end{itemize}

\textbf{Structure de \texttt{EnvironmentExp}} :

Liste chaînée de dictionnaires (\texttt{Map<Symbol, ExpDefinition>}). Chaque bloc introduit un nouvel environnement fils. La recherche remonte vers le parent si non trouvé.

\textbf{Mécanisme de décoration} :

Lors de \texttt{verify*()}, chaque expression se voit attribuer un type via \texttt{setType()}. Ces types sont utilisés en génération de code.

\subsection{\texttt{fr.ensimag.deca.codegen}}

Contient uniquement \texttt{RegisterManager} pour l'allocation de registres.

\textbf{Dépendances externes} :

\texttt{fr.ensimag.ima.pseudocode} fournit les instructions IMA (\texttt{LOAD}, \texttt{STORE}, \texttt{ADD}, \texttt{BEQ}, etc.).

\subsection{\texttt{fr.ensimag.ima.pseudocode}}

Représentation abstraite de l'assembleur IMA.

\begin{itemize}[itemsep=0.2em]
    \item \texttt{IMAProgram} : Liste de lignes (instructions + labels + commentaires)
    \item \texttt{Instruction} : Classes comme \texttt{LOAD}, \texttt{ADD}, \texttt{BEQ}, etc.
    \item \texttt{GPRegister} : Registres généraux R0-R15
    \item \texttt{DAddr} : Adresses (registres, offsets)
\end{itemize}

\section{Algorithmes et Structures de Données Spécifiques}

\subsection{Allocation de Registres : Stratégie Naïve}

\textbf{Choix} : Allocation naïve avec spill automatique sur pile.

\textbf{Implémentation (\texttt{RegisterManager})} :

\begin{itemize}[itemsep=0.2em]
    \item Registres R2 à R15 disponibles (R0-R1 réservés scratch)
    \item Compteur \texttt{registreCourant} : prochain registre libre
    \item Méthodes :
    \begin{itemize}
        \item \texttt{prendreRegistre()} : alloue un registre, échoue si tous utilisés
        \item \texttt{libererRegistre()} : libère le dernier registre alloué (stack-like)
    \end{itemize}
\end{itemize}

\textbf{Stratégie dans \texttt{AbstractOpArith}} :

\begin{enumerate}[itemsep=0.3em]
    \item Évaluer opérande gauche dans \texttt{register}
    \item Si registre libre : évaluer opérande droite dans nouveau registre, puis \texttt{OP reg\_right, register}
    \item Si aucun registre libre (spill) :
    \begin{itemize}
        \item \texttt{PUSH register} (sauvegarder gauche sur pile)
        \item Évaluer opérande droite dans \texttt{register}
        \item \texttt{LOAD register, R0} puis \texttt{POP register}
        \item \texttt{OP R0, register}
    \end{itemize}
\end{enumerate}

\textbf{Justification} :

\begin{itemize}[itemsep=0.2em]
    \item \textbf{Simplicité} : Pas de graphe d'interférence, pas d'analyse de vivacité
    \item \textbf{Suffisant pour Deca} : Expressions simples, peu d'imbrications profondes
    \item \textbf{Inconvénient} : Spill fréquent pour expressions complexes (overhead pile)
    \item \textbf{Alternative rejetée} : Allocation par coloration de graphe (complexité excessive pour GL)
\end{itemize}

\subsection{Génération de Code pour le Flot de Contrôle}

\textbf{Stratégie Labels} :

Les structures \texttt{if} et \texttt{while} utilisent des labels uniques pour les sauts.

\textbf{\texttt{IfThenElse}} :

\begin{lstlisting}[style=bashstyle, numbers=none]
<Code(Condition, faux, E_Sinon)>
<Code(Alors)>
BRA E_Fin
E_Sinon:
<Code(Sinon)>
E_Fin:
\end{lstlisting}

\begin{itemize}[itemsep=0.2em]
    \item Compteur statique \texttt{c} incrémenté pour labels uniques (\texttt{else.1}, \texttt{end\_if.1}, etc.)
    \item \texttt{codeGenBool(compiler, false, elseLabel)} : saute si condition fausse
\end{itemize}

\textbf{\texttt{While}} :

\begin{lstlisting}[style=bashstyle, numbers=none]
BRA E_Cond.n
E_Debut.n:
<Code(Corps)>
E_Cond.n:
<Code(Condition, vrai, E_Debut.n)>
\end{lstlisting}

\begin{itemize}[itemsep=0.2em]
    \item Évaluation condition en fin de boucle (évite duplication)
    \item Saut initial vers condition (cohérence sémantique)
\end{itemize}

\textbf{Génération de Conditions} :

\texttt{codeGenBool(compiler, branchOn, targetLabel)} évalue une expression booléenne :

\begin{enumerate}[itemsep=0.2em]
    \item Évaluer expression dans registre
    \item \texttt{CMP \#0, register}
    \item Si \texttt{branchOn == true} : \texttt{BNE targetLabel} (sauter si vrai)
    \item Si \texttt{branchOn == false} : \texttt{BEQ targetLabel} (sauter si faux)
\end{enumerate}

\textbf{Justification} :

\begin{itemize}[itemsep=0.2em]
    \item \textbf{Labels uniques} : Évite conflits dans boucles imbriquées
    \item \textbf{Évaluation paresseuse non implémentée} : \&\& et || évaluent toujours les deux opérandes (simplification)
    \item \textbf{Alternative} : Court-circuit (complexe, non requis par sujet)
\end{itemize}

\subsection{Gestion de la Pile : TSTO/ADDSP}

\textbf{Stratégie} :

\texttt{RegisterManager} trace la taille maximale de pile nécessaire.

\textbf{Compteurs} :

\begin{itemize}[itemsep=0.2em]
    \item \texttt{taillePileCourante} : Profondeur actuelle de pile
    \item \texttt{taillePileMax} : Maximum atteint (pour TSTO)
    \item \texttt{nbGlobales} : Nombre de variables globales
\end{itemize}

\textbf{Méthodes} :

\begin{itemize}[itemsep=0.2em]
    \item \texttt{empiler()} : Incrémente compteurs (lors de PUSH ou spill)
    \item \texttt{depiler()} : Décrémente (lors de POP)
    \item \texttt{ajouterVariablesLocales(n)} : Réserve espace pour variables locales
\end{itemize}

\textbf{Génération en-tête (\texttt{Program.codeGenProgram()})} :

Instructions insérées en \textbf{ordre inverse} (LIFO via \texttt{addFirstInstruction()}) :

\begin{enumerate}[itemsep=0.2em]
    \item \texttt{TSTO \#(maxTemp + nbGlob)} : Vérifie espace disponible
    \item \texttt{BOV stack\_overflow\_error} : Saute vers gestion erreur si overflow
    \item \texttt{ADDSP \#nbGlob} : Réserve espace pour variables globales
\end{enumerate}

\textbf{Gestion d'erreurs} :

Labels pour erreurs runtime :

\begin{lstlisting}[style=bashstyle, numbers=none]
stack_overflow_error:
    WSTR "Error: Stack Overflow"
    WNL
    ERROR
\end{lstlisting}

\textbf{Justification} :

\begin{itemize}[itemsep=0.2em]
    \item \textbf{Analyse statique} : Calcul conservateur de taille pile (pas de récursion donc borné)
    \item \textbf{Sécurité} : \texttt{BOV} empêche corruption mémoire
    \item \textbf{Coût} : Une instruction TSTO par programme (négligeable)
    \item \textbf{Alternative rejetée} : Pile dynamique (non supporté par IMA)
\end{itemize}

\subsection{Conversion de Types Implicites}

\textbf{Stratégie} :

Si opération arithmétique mélange \texttt{int} et \texttt{float}, convertir \texttt{int} en \texttt{float}.

\textbf{Implémentation (\texttt{AbstractOpArith.verifyExpr()})} :

\begin{lstlisting}[style=javastyle]
if (t1.isFloat() && t2.isInt()) {
    setRightOperand(new ConvFloat(getRightOperand()));
    // Reverifier l'operande convertie
}
\end{lstlisting}

Node \texttt{ConvFloat} génère instruction IMA \texttt{FLOAT register, register}.

\textbf{Justification} :

\begin{itemize}[itemsep=0.2em]
    \item \textbf{Norme Deca} : Promotions implicites int→float (comme Java)
    \item \textbf{Simplicité AST} : Insertion transparente de nœuds de conversion
    \item \textbf{Alternative} : Conversion runtime (moins efficace, nécessite métadonnées type)
\end{itemize}

\section{Stratégie de Développement}

\subsection{Approche Incrémentale}

Le projet suit un développement par étapes :

\begin{enumerate}[itemsep=0.3em]
    \item \textbf{Étape A (Lexer/Parser)} : Grammaires ANTLR4, tests lexicaux/syntaxiques
    \item \textbf{Étape C "Sans Objet"} : Génération de code procédural (sans classes)
    \item \textbf{Étape B (Analyse Contextuelle)} : Vérification types et portées
    \item \textit{Étape C Complète (Non implémentée)} : Support OO (classes, héritage, méthodes virtuelles)
\end{enumerate}

\subsection{Stratégie de Tests}

\textbf{Tests par phase} :

\begin{itemize}[itemsep=0.3em]
    \item \textbf{Lexicaux} : \texttt{src/test/script/basic-lex.sh} + \texttt{test\_lex} launcher
    \begin{itemize}
        \item Fichiers valides dans \texttt{src/test/deca/lexical/valid/}
        \item Fichiers invalides dans \texttt{src/test/deca/lexical/invalid/}
    \end{itemize}
    \item \textbf{Syntaxiques} : \texttt{basic-synt.sh} + \texttt{test\_synt}
    \item \textbf{Contextuels} : \texttt{basic-context.sh} + \texttt{test\_context}
    \begin{itemize}
        \item Vérifie erreurs de types, variables non déclarées
    \end{itemize}
    \item \textbf{Génération} : \texttt{basic-gencode.sh}
    \begin{itemize}
        \item Compile et exécute avec IMA, vérifie sortie
    \end{itemize}
\end{itemize}

\textbf{Tests unitaires JUnit 5} :

Fichiers dans \texttt{src/test/java/fr/ensimag/deca/} :

\begin{itemize}[itemsep=0.2em]
    \item \texttt{TestWhileCodeGen} : Teste génération de boucles
    \item \texttt{TestBooleanOpsCodeGen} : Opérateurs logiques
    \item \texttt{TestPlusPlain}, \texttt{TestPlusAdvanced} : Vérification contextuelle
\end{itemize}

\textbf{Test de régression} :

\texttt{./src/test/script/common-tests.sh} : Suite minimale validant fonctionnalités de base (hello world, syntaxe invalide, variables non déclarées).

\subsection{Couverture de Code}

JaCoCo configuré via Maven (\texttt{jacoco.skip=false}). Rapport généré dans \texttt{target/site/jacoco/}.

\textbf{Exécution} :

\begin{lstlisting}[style=bashstyle]
mvn clean test -Djacoco.skip=false
./src/test/script/jacoco-report.sh
\end{lstlisting}

\textbf{Exclusions} : Classes d'erreurs internes (\texttt{DecacInternalError}, \texttt{IMAInternalError}) exclues de la couverture.

\section{Spécifications Techniques}

\subsection{Contraintes d'Environnement}

\begin{itemize}[itemsep=0.2em]
    \item \textbf{Java} : Version 21 requise (déclaré dans \texttt{pom.xml})
    \item \textbf{Maven} : 3.6.3 minimum (vérifié par \texttt{maven-enforcer-plugin})
    \item \textbf{ANTLR4} : 4.13.2 (dépendance Maven)
    \item \textbf{JUnit} : 5 (Jupiter) pour tests unitaires
\end{itemize}

\subsection{Conventions de Codage}

\textbf{Nommage} :

\begin{itemize}[itemsep=0.2em]
    \item Classes abstraites préfixées \texttt{Abstract} (ex: \texttt{AbstractExpr})
    \item Méthodes \texttt{verify*()} retournent \texttt{Type} et lèvent \texttt{ContextualError}
    \item Méthodes \texttt{codeGen*()} sont \texttt{protected} et ne retournent rien
\end{itemize}

\textbf{Gestion d'erreurs} :

\begin{itemize}[itemsep=0.2em]
    \item \textbf{Erreurs lexicales/syntaxiques} : ANTLR rapporte sur \texttt{System.err} (format \texttt{file:line:message})
    \item \textbf{Erreurs contextuelles} : \texttt{throw new ContextualError(msg, location)}
    \item \textbf{Erreurs fatales} : \texttt{throw new DecacFatalError(msg)} (I/O, assertions violées)
\end{itemize}

\textbf{Assertions Java} :

Le code utilise massivement \texttt{assert}. JVM doit être lancée avec \texttt{-ea} (activé dans script \texttt{decac} et launchers de test).

\textbf{Pattern Visitor} :

Bien que non explicite, l'AST suit un pattern Visitor implicite :

\begin{itemize}[itemsep=0.2em]
    \item \texttt{verifyExpr()} traverse l'arbre en pré-ordre
    \item \texttt{codeGenExpr()} traverse en post-ordre (évaluation des fils avant le parent)
\end{itemize}

\subsection{Limitations Connues}

\begin{itemize}[itemsep=0.2em]
    \item \textbf{Pas de compilation parallèle} : Flag \texttt{-P} lève \texttt{UnsupportedOperationException}
    \item \textbf{OOP partiel} : Classes et héritage non complètement implémentés (étape C objet manquante)
    \item \textbf{Pas d'optimisation} : Code généré non optimisé (pas de pliage de constantes, pas d'élimination de code mort)
    \item \textbf{Registres limités} : Spill naïf peut générer beaucoup de PUSH/POP pour expressions complexes
\end{itemize}

\subsection{Fichiers de Configuration}

\begin{itemize}[itemsep=0.2em]
    \item \texttt{pom.xml} : Configuration Maven (dépendances, plugins ANTLR4, JaCoCo)
    \item \texttt{src/main/resources/log4j.properties} : Configuration logging
    \item \texttt{src/main/config/findbugs-exclude.xml} : Exclusions analyse statique
    \item \texttt{src/test/script/launchers/*} : Scripts configurant classpath pour tests par phase
\end{itemize}

\section{Conclusion}

Cette architecture privilégie la \textbf{simplicité et la maintenabilité} :

\begin{itemize}[itemsep=0.2em]
    \item Pipeline classique de compilation
    \item Allocation naïve de registres (sans graphe d'interférence)
    \item Tests modulaires par phase
    \item Séparation claire AST / Context / CodeGen
\end{itemize}

Le compilateur remplit son objectif pédagogique : traduire Deca procédural vers IMA. Les extensions futures (OOP complet, optimisations) nécessiteront :

\begin{itemize}[itemsep=0.2em]
    \item Refonte de \texttt{RegisterManager} pour allocation globale
    \item Ajout de passes d'optimisation (DCE, CSE, inlining)
    \item Support tables de méthodes virtuelles (VMT)
\end{itemize}

\vspace{1cm}
\noindent\textbf{Contacts} : Équipe GL 51, Ensimag 2026

\end{document}
